\color{unimain}{\subsection{Фиксация положения выключателя (ФПВ)}}
\color{black}

\begin{enumerate}[label=\arabic{section}.\arabic{subsection}.\arabic{enumi}, labelsep=4pt, leftmargin=0pt, itemindent=57pt, itemsep=0pt, parsep=0pt] % если что можно крутить itemsep 0 или 5

\item
Функциональный блок фиксации положения выключателей предназначен для определения состояния
выключателей на основании технологической информации:
\begin{itemize}
\item состояния блок-контактов выключателя;
\item положения разъединителей в цепи выключателя;
\item информации о выводе в ремонт выключателя.
\end{itemize}

\item
Ввод и вывод функции ФПВ осуществляется либо посредством программного ключа \textbf{«ФПВ Вn»}, либо с
использованием внешнего сигнала \textit{«Вывод ФПВ Вn»}.

\item
Контроль положения разъединителей осуществляется посредством программного ключа 
\textbf{«Контр. Р Вn»}.


\item
Предусмотрена возможность приема сигналов положения блок-контактов как для трехфазного, так и для
пофазного исполнения коммутационных аппаратов.

\vspace{2mm}
\begin{figure}[H]
\centering
\includegraphics[width=0.7\textwidth,height=0.7\textheight,keepaspectratio]{\fbpath/762.3461 ФПВ/img/FPV.pdf}
\captionof{figure}{Функциональная схема ФПВ}\label{pict_fpv:FPV}
\end{figure}

\small
\begin{longtable}{|>{\centering\arraybackslash}m{5.3cm}|>{\centering\arraybackslash}m{3.3cm}|>{\centering\arraybackslash}m{4.2cm}|>{\centering\arraybackslash}m{1.8cm}|>{\centering\arraybackslash}m{1cm}|}
\caption{Параметры для настройки ФПВ\hfill\vspace{-0.5\baselineskip}}\label{fpv:tbl1}\\ 
\hline
\rowcolor{gray!30}
Параметр & Условное обозначение на схеме & Значение/ Диапазон & Единица измерения & Шаг \\ 
\hline
\endfirsthead
\caption*{\hspace{3pt}\emph{Продолжение таблицы \ref{fpv:tbl1}\hfill\vspace{-0.5\baselineskip}}} \\ % сделано по ГОСТ 2.105 п.6.8.7
\hline
\rowcolor{gray!30}
Параметр & Условное обозначение на схеме & Значение/ Диапазон & Единица измерения & Шаг \\ 
%\hline
\endhead

%===m>ФПВ|-
%===m
\end{longtable}
\normalsize

\end{enumerate}

\newpage